% !TeX root = main.tex

An important features of OSTRICH is the possibility to extend the
solver with minimal effort.
%
OSTRICH’s \texttt{PreOp} trait lets users add new string functions by
overriding three core methods. This can be done in either Scala (the
main language used to implement OSTRICH) or in Java.
%
Below we show each method in turn,
together with the \texttt{ReversePreOp} implementation\footnote{This
  snippet is abbreviated for clarity. The full implementation in
  \texttt{ReversePreOp.scala} includes additional automaton‐type
  matching and error handling.}.

First, the method \texttt{apply} computes the \emph{pre-image} of a
regular language under $f$, which is the core operation needed for
backward propagation in RCP.  The method is passed any existing automata constraints on the $r$ arguments (which may be ignored) and the automaton $\mathcal{A}$ representing the result language.  The method returns an iterator over tuples of automata whose images under $f$ lie inside $\mathcal{A}$, plus (optionally) the subset of the input constraints actually used.

\begin{lstlisting}[language=Scala]
override def apply(
	argCs: Seq[Seq[Automaton]],
	resultC: Automaton
) : (Iterator[Seq[Automaton]], Seq[Seq[Automaton]]) = {
	(Iterator(Seq(ReverseAutomaton(resultC))), List())
}
\end{lstlisting}

Second, the method \texttt{eval} performs \emph{concrete} evaluation on ground strings: when all arguments are known, it simply returns the result of applying $f$.

\begin{lstlisting}[language=Scala]
override def eval(
	args: Seq[Seq[Int]]
): Option[Seq[Int]] = Some(args.head.reverse)
\end{lstlisting}

Third, \texttt{forwardApprox} computes a (sound) \emph{post-image} of input languages under $f$, used for forward propagation.  While returning the universal automaton is always correct, a tighter approximation greatly improves performance.  For \texttt{str.reverse}, one can exactly compute this by intersecting any input automata and then reversing the resulting automaton.

\begin{lstlisting}[language=Scala]
override def forwardApprox(
	argCs: Seq[Seq[Automaton]]
): Automaton = {
	val prod = ProductAutomaton(argCs)
	ReverseAutomaton(prod)
}
\end{lstlisting}
%
%Putting it all together:
%
%\begin{lstlisting}[language=Scala]
%package ostrich.preop
%
%object ReversePreOp extends PreOp {
%	override def name: String  = "str.reverse"
%	override def arity: Int     = 1
%
%	// 1) Pre-image
%	override def apply(
%		argCs: Seq[Seq[Automaton]],
%		resultC: Automaton
%	): (Iterator[Seq[Automaton]], Seq[Seq[Automaton]]) = {
%		(Iterator(Seq(ReverseAutomaton(resultC))), List())
%	}
%
%	// 2) Concrete evaluation
%	override def eval(
%		args: Seq[Seq[Int]]
%	): Option[Seq[Int]] = Some(args.head.reverse)
%
%	// 3) Post-image
%	override def forwardApprox(
%		argCs: Seq[Seq[Automaton]]
%	): Automaton = {
%		val prod = ProductAutomaton(argCs)
%		ReverseAutomaton(prod)
%	}
%}
%\end{lstlisting}


Finally, the new string function has to be registered in the core
string theory class of OSTRICH.
In the file \texttt{OstrichStringTheory.scala}, we register \texttt{str.reverse} by first declaring its SMT-side function symbol:

\begin{lstlisting}[language=Scala]
val str_reverse = MonoSortedIFunction(
	"str.reverse",
	List(StringSort), StringSort, true, false
)
\end{lstlisting}

We then add it to the \texttt{extraStringFunctions} list, pairing the name, the \texttt{IFunction}, our \texttt{ReversePreOp} implementation, and the simple lambdas that pick out the input and output terms:

\begin{lstlisting}[language=Scala]
val extraStringFunctions = List(
	("str.reverse", str_reverse, ReversePreOp,a =>
		 List(a(0)),a => a(1)))
\end{lstlisting}

With these two lines in place, any occurrence of \texttt{str.reverse} in the SMT-LIB script is recognized by OSTRICH2 and automatically dispatched to our \texttt{ReversePreOp} for pre-image computation, concrete evaluation, and forward-approximation.


